\section*{How to use this mini-textbook}
\addcontentsline{toc}{section}{How to use this mini-textbook}
This is a \emph{teaching document} for the full PHYS4055
course, not an abbreviated formula sheet.
Its purpose is to answer three questions:
what is happening physically, why the formula
has its particular form, and how that formula
helps solve an exam problem.
The 15-page overview is designed for fast lookup;
this book is intended to be \emph{read and worked through}.

\subsection*{A suggested route through the material}
The course has four connected parts.
Quantum optics first describes light as quantised
field modes. Applications of optics use those
states, operations and measurements to
communicate, process information and make
sensitive measurements. Atomic physics
introduces electron wavefunctions, angular
momentum, energy-level splitting and transition
matrix elements. Applications of atomic physics
use these ideas in cooling and trapping,
symmetry tests and clocks.

If quantum notation feels unfamiliar, start
with Chapters 1--4 before trying the
applications. If your main difficulty is
angular-momentum coupling, read Chapters
14--16 slowly, including their worked
examples. Later chapters repeatedly
refer back to these ideas.

\subsection*{Mathematical notation, explained}
A ket \(\ket\psi\) is a quantum state vector.
A bra \(\bra\phi\) is its conjugate transpose.
The inner product \(\braket\phi\psi\)
is a complex amplitude and
\(|\braket\phi\psi|^2\) is a probability
when the states are normalised.
A matrix element \(\mel fO i\) asks
how strongly the operator \(O\)
connects initial \(\ket i\) to final
\(\ket f\). A commutator is
\([A,B]=AB-BA\); it measures
the difference caused by operator order.

A density operator \(\rho\) describes
both pure states and mixtures.
An expectation is
\(\langle O\rangle=\operatorname{Tr}(\rho O)\).
A partial trace sums over outcomes
of a subsystem that is not observed.
A unitary \(U^\dagger U=I\)
preserves state norms.
A Hermitian observable \(O=O^\dagger\)
has real measurement eigenvalues.

\begin{worked}{How to read a transition matrix element}
The expression
\(\bra{3s}\mathbf d\ket{3p}\)
does not directly give a transition
probability. It gives a \emph{dipole
amplitude}, including radial overlap
and angular selection rules.
Its squared magnitude enters an
Einstein \(A\) rate, multiplied by
a factor proportional to
the cube of transition frequency.
Only after summing all decay
channels can we infer a lifetime.
\end{worked}

\subsection*{Conventions we maintain}
For the quantum-optics chapters,
\[
x=a+a^\dagger,\qquad
p=i(a^\dagger-a),\qquad
[x,p]=2i,
\]
so vacuum quadrature variance is 1.
Optical beam-splitter phases are
defined locally in each example;
physical probabilities are invariant
under a consistent convention change.

For atomic physics, lower-case \(l,s\)
often refer to one electron, capital
\(L,S\) to the combined electrons,
\(J\) to electronic total angular
momentum, \(I\) to nuclear spin and
\(F\) to their sum. The two symbols
\(m_J,m_F\) name projections along
the chosen quantisation axis.

\subsection*{The four checks to do at the end}
\begin{enumerate}
\item \textbf{Probability:} do all output
probabilities sum to 1? Is \(\operatorname{Tr}\rho=1\)?
\item \textbf{Symmetry:} is the transition
allowed by triangle, magnetic projection
and parity conditions?
\item \textbf{Physical limiting case:}
what happens if the loss is zero,
the detuning is zero, or the field is removed?
\item \textbf{Units:} do the answers have the
expected dimensions and order of magnitude?
\end{enumerate}
Each chapter ends with a short concept
check. For longer algebraic practice,
use the companion Worked Examples,
Assignment Solutions and Past Exam
Solutions booklets in the repository.
